%% A281_Traeger_und_Information_De.tex
%% FFGFT A-Serie -- A281: Träger und Information
%% Autor: Johann Pascher  --  ORCID 0009-0000-6518-4064
%% Sprache: Deutsch  --  Engine: LuaLaTeX
%% Vorlage: landauer_erweitert.tex (undatiert)
%% Begleitdokument zu A280

\documentclass[11pt,a4paper]{article}

% ---- Engine: LuaLaTeX ----
\usepackage{fontspec}
\usepackage{mathtools}
\usepackage{unicode-math}
\usepackage{polyglossia}
\setmainlanguage{german}

\usepackage{microtype}
\emergencystretch=3.5em
\tolerance=2500
\hbadness=3000
\usepackage{amsthm}
\usepackage{enumitem}
\usepackage{booktabs}
\usepackage{longtable}
\usepackage{array}
\usepackage[a4paper,margin=2.6cm,headsep=1.1em,headheight=22pt]{geometry}
\usepackage{fancyhdr}
\usepackage{titlesec}
\usepackage{xcolor}
\usepackage{needspace}
\usepackage{tocloft}

% ---- Fonts ----
\setmainfont{Inter}[
  Scale=1.02,
  Ligatures=TeX ]
\setsansfont{Inter}[Scale=MatchLowercase,Ligatures=TeX]
\setmonofont{JetBrains Mono}[Scale=0.88]
\setmathfont{Libertinus Math}

% ---- Colours ----
\definecolor{accent}{HTML}{1F3A5F}
\definecolor{accent2}{HTML}{6A4E2E}
\definecolor{rulecol}{HTML}{B9C2CE}
\definecolor{linknav}{HTML}{2C5F8A}
\definecolor{linkcite}{HTML}{7A5C2E}

% ---- Schichtmarkierungen (A-Serie) ----
\definecolor{laySet}{HTML}{7A2E2E}
\definecolor{layK}{HTML}{1D5C3A}
\definecolor{layB}{HTML}{1F3A5F}
\definecolor{layS}{HTML}{6A4E2E}
\definecolor{layH}{HTML}{5A2E6A}

\newcommand{\LS}{\textcolor{laySet}{\textbf{\footnotesize\textsf{[SETZUNG]}}}\,}
\newcommand{\LK}{\textcolor{layK}{\textbf{\footnotesize\textsf{[K]}}}\,}
\newcommand{\LB}{\textcolor{layB}{\textbf{\footnotesize\textsf{[B]}}}\,}
\newcommand{\LSk}{\textcolor{layS}{\textbf{\footnotesize\textsf{[S]}}}\,}
\newcommand{\LH}{\textcolor{layH}{\textbf{\footnotesize\textsf{[H]}}}\,}

% ---- Theorem-Umgebungen ----
\newtheoremstyle{ffgft}{0.8em}{0.4em}{\itshape}{}{\bfseries\color{accent}}{.}{0.5em}{}
\newcommand{\qedsq}{\raisebox{0.03em}{\framebox[0.62em]{\rule{0pt}{0.42em}}}}
\theoremstyle{ffgft}
\newtheorem{definition}{Definition}
\newtheorem{satz}{Satz}

% ---- Hyperref ----
\usepackage[
  colorlinks=true,
  linkcolor=linknav,
  citecolor=linkcite,
  urlcolor=linknav,
  bookmarksnumbered=true,
  bookmarksopen=true,
  pdfstartview=FitH ]{hyperref}
\hypersetup{
  pdftitle={Träger und Information (A281)},
  pdfauthor={Johann Pascher},
  pdfsubject={Landauer-Prinzip --- Textkritik und Rezeptionsgeschichte},
  pdfkeywords={Landauer, Träger, Information, Hypostasierung, Rezeption, FFGFT, A281}
}

% ---- Section styling ----
\titleformat{\section}
  {\color{accent}\normalfont\Large\bfseries}
  {\color{accent2}\thesection}{0.7em}{}
\titleformat{\subsection}
  {\color{accent}\normalfont\large\bfseries}
  {\color{accent2}\thesubsection}{0.6em}{}
\titlespacing*{\section}{0pt}{1.6\baselineskip}{0.7\baselineskip}
\titlespacing*{\subsection}{0pt}{1.1\baselineskip}{0.4\baselineskip}

% ---- Headers ----
\pagestyle{fancy}
\fancyhf{}
\renewcommand{\headrulewidth}{0.4pt}
\renewcommand{\headrule}{\hbox to\headwidth{\color{rulecol}\leaders\hrule height \headrulewidth\hfill}}
\fancyhead[L]{\footnotesize\color{accent}\itshape FFGFT $\cdot$ A281 $\cdot$ Träger und Information}
\fancyhead[R]{\footnotesize\color{accent2}\thepage}

% ---- TOC look ----
\renewcommand{\cftsecfont}{\normalfont}
\renewcommand{\cftsecpagefont}{\normalfont}
\renewcommand{\cftsecleader}{\cftdotfill{\cftdotsep}}
\setlength{\cftbeforesecskip}{3pt}

\setlength{\parskip}{0.35em}
\setlength{\parindent}{0pt}
\linespread{1.06}

\begin{document}

% ===================== TITEL =====================
\begin{titlepage}
\centering
\vspace*{2.2cm}
{\color{rulecol}\rule{0.5\textwidth}{0.6pt}}\\[1.4cm]
{\color{accent}\Huge\bfseries Träger und Information\par}
\vspace{1.0cm}
{\large\itshape Was Landauer 1961 beweist\\ --- und was die Rezeption daraus gemacht hat\par}
\vspace{0.6cm}
{\color{accent2}\rule{0.28\textwidth}{0.4pt}}\\[1.0cm]
{\large FFGFT A-Serie $\cdot$ A281\par}
\vspace{0.4cm}
{\normalsize Begleitdokument zu A280 --- Textkritik, Rezeptionsgeschichte\\ und Sprachkritik\par}
\vfill
{\large Johann Pascher\par}
\vspace{0.3cm}
{\footnotesize ORCID: 0009-0000-6518-4064\par}
\vspace{0.2cm}
{\footnotesize Juli 2026\par}
\vspace{1.4cm}
{\color{rulecol}\rule{0.5\textwidth}{0.6pt}}
\vspace*{1.5cm}
\end{titlepage}

% ===================== TOC =====================
\tableofcontents
\thispagestyle{fancy}
\clearpage

% ===================== KURZFASSUNG =====================
\begingroup\small\noindent
\textbf{Kurzfassung.} Die geläufige Aussage „das Löschen eines Bits erzeugt
mindestens $kT\ln 2$ Wärme“ geht über das hinaus, was Landauers Originalarbeit von
1961 beweist. Landauer zeigt präzise: Wird ein physikalisches
Zwei-Zustands-System im thermischen Gleichgewicht unabhängig von seinem
Ausgangszustand in einen definierten Zustand gezwungen, muss mindestens
$kT\ln 2$ dissipiert werden. Das ist eine Aussage über \emph{Träger-Operationen}.
Der Schritt von dort zur universellen Behauptung über \emph{Information} als
solche wird in Landauers Arbeit nicht vollzogen; er wird durch eine begriffliche
Doppeldeutigkeit des Wortes „bit“ begünstigt, ist aber keine bewiesene Konsequenz
seiner Argumente. Diese Arbeit trennt beides, weist Landauers eigene
Einschränkung nach, verfolgt die Verallgemeinerung durch die Rezeption und
benennt den sprachlichen Mechanismus --- die Hypostasierung einer Abstraktion ---,
der sie getragen hat. \emph{Verhältnis zu A280:} A281 fügt der Physik nichts hinzu.
A280 führt die Buchhaltung; A281 führt die Textkritik. Die Wörterbuch-Zuordnung
beider Fassungen steht in Abschnitt~\ref{sec:woerterbuch}.
\par\endgroup

\bigskip

\begingroup\small\noindent
\textbf{Schichtmarkierungen.} \LS deklarierter Ausgangspunkt ---
\LB algebraisch/mathematisch aus deklarierten Grundlagen abgeleitet ---
\LK empirisch geprüft bzw.\ aus Primärquellen belegt ---
\LSk Plausibilitätsskizze ---
\LH offene Forschungsfrage.
\par\endgroup

\bigskip

% ===================== 1 =====================
\section{Die Fragestellung}

Seit Landauers Arbeit von 1961 [B1] wird in der Literatur häufig behauptet, jedes
gelöschte Bit müsse mindestens $kT\ln 2$ an Wärme erzeugen. Die Aussage ist fester
Bestandteil von Informatik, Physik und Wissenschaftsphilosophie geworden.

Die vorliegende Arbeit vertritt: Diese Verallgemeinerung geht über das hinaus, was
Landauers Beweis zeigt. Landauer macht eine präzise Aussage über physikalische
Träger-Operationen. Der Sprung von dort zur universellen Aussage über Information
als abstrakte Größe ist nicht sein Schritt --- er wurde in der Rezeption
vollzogen, begünstigt durch eine begriffliche Unschärfe, die in der Disziplin
angelegt ist.

Zwei Abgrenzungen vorweg, damit die Reichweite der Arbeit nicht überschätzt wird.

\LB \textbf{Erstens: Dies ist keine Widerlegung des Landauer-Prinzips.} Was hier
bestritten wird, ist die \emph{populäre Kurzfassung}, nicht Landauers Rechnung.
Für die physikalische Operation, die Landauer betrachtet, ist sein Ergebnis
korrekt. Bestritten wird die Zuständigkeit dieses Ergebnisses für einen Gegenstand
--- die abstrakte Information ---, über den es keine Aussage macht.

\LB \textbf{Zweitens: Die physikalische Substanz steht in A280, nicht hier.} Dass
die thermodynamische Untergrenze am physikalischen Gebiet hängt und nicht an
seiner Beschreibung, ist dort aus Liouville hergeleitet. A281 fügt dem nichts
hinzu. Was A281 leistet, ist anderes und eigenständig: die Textanalyse von
Landauer 1961, der Nachweis von Landauers eigener Einschränkung, die
Rezeptionsgeschichte der Verallgemeinerung und die Sprachkritik der Digitaltechnik.

\subsection{Wörterbuch: A280 und A281}
\label{sec:woerterbuch}

\LS Die beiden Dokumente treffen dieselbe Unterscheidung in verschiedenen Wörtern.
Damit im Korpus keine zwei parallelen Terminologien für einen Sachverhalt
entstehen, wird die Zuordnung hier festgelegt und ist verbindlich:

\begin{center}
\small
\begin{tabular}{@{}p{0.36\textwidth}p{0.36\textwidth}p{0.20\textwidth}@{}}
\toprule
\textbf{A280 (Buchhaltung)} & \textbf{A281 (Textkritik)} & \textbf{Ebene} \\
\midrule
Gebiet im Phasenraum & Träger & physikalisch \\
Beschreibung & Information & abstrakt \\
Gebietsreduktion $\{A,B\}\to Z$ & Träger-Operation \textup{RESTORE TO ONE} & physikalisch \\
Nicht-Bijektivität der Gebietsabbildung & Verkleinerung des Träger-Zustandsraums & physikalisch \\
Beschreibungsentscheidung & interpretative Löschung & abstrakt \\
\bottomrule
\end{tabular}
\end{center}

\LB Die Zuordnung ist eine Übersetzung, keine Erweiterung. Wer „Träger“ liest,
darf „Gebiet“ denken, und umgekehrt. Eine Aussage, die in einer Spalte gilt, gilt
in der anderen --- andernfalls liegt ein Fehler in einem der beiden Dokumente vor,
nicht ein neuer Sachverhalt.

% ===================== 2 =====================
\section{Die Kernunterscheidung}

\subsection{Definitionen}

\begin{definition}[Information]
Information ist eine abstrakte Größe. Ein Bit ist eine Einheit der Information,
die einen von zwei Zuständen ($0$ oder $1$) repräsentiert. Information hat keine
Masse, keine Energie, keine Entropie und keine physikalische Lokalisation.
\end{definition}

\begin{definition}[Träger]
Der Träger der Information ist ein physikalisches System, das die Information
repräsentiert: ein elektrischer Schaltkreis, ein magnetischer Spin, ein Photon
oder jedes andere physikalische Zwei-Zustands-System. Der Träger hat Energie,
Entropie und weitere physikalische Eigenschaften.
\end{definition}

\LS Beide Definitionen sind Setzungen. Sie werden nicht bewiesen; sie legen fest,
worüber im Folgenden geredet wird. Der gesamte Ertrag der Arbeit hängt daran, dass
diese Trennung durchgehalten wird.

\subsection{Energie hängt am Träger}

\begin{satz}
Einem Informationsbit kann keine Energie zugeordnet werden, nur seinem Träger.
\end{satz}

\LB \textbf{Beweis.} Angenommen, ein Informationsbit hätte eine intrinsische
Energie. Dann müsste diese Energie unabhängig von der physikalischen Realisierung
sein. Ein Bit mit dem Wert $1$ lässt sich jedoch durch verschiedene physikalische
Systeme mit verschiedenen Energien repräsentieren: durch 5\,V in einem
CMOS-Schaltkreis, durch 1\,V in einer anderen Technologie, durch einen
magnetischen Spin mit einer bestimmten Austauschenergie, durch ein Photon einer
bestimmten Frequenz. Da der Informationsgehalt in allen Fällen identisch ist, die
Energie aber variiert, kann die Energie keine Eigenschaft der Information sein.
Sie ist eine Eigenschaft des Trägers. \hfill\qedsq

Das Argument ist ein Mehrfachrealisierungs-Argument und trägt genau so weit, wie
Mehrfachrealisierung reicht --- also vollständig: Es gibt keine
Realisierungsklasse, in der die Energie mit dem Informationsgehalt mitvariieren
würde.

\subsection{Das Bus-Beispiel}

\LB Ein anschauliches Beispiel ist die Bitbreite eines Datenbusses:

\begin{itemize}[leftmargin=1.4em,itemsep=0.25em]
  \item Die \textbf{Busbreite} (etwa 64 Bit) ist eine feste physikalische
        Eigenschaft der Hardware. Sie bestimmt, wie viele Träger gleichzeitig
        umgeschaltet werden.
  \item Die \textbf{Information} ist das, was durch die Muster kodiert wird.
        Dasselbe Bitmuster kann je nach Kontext Verschiedenes bedeuten.
\end{itemize}

Ein 64-Bit-Bus schaltet immer 64 Träger, unabhängig davon, ob die Information eine
Zahl, ein Buchstabe oder ein Bildpunkt ist. Die Leitungen haben eine bestimmte
Kapazität und brauchen Energie für die Umladung. Die Information hat diese Energie
nicht.

% ===================== 3 =====================
\section{Landauers Modell --- was es zeigt}

\subsection{Das Teilchen im Doppeltopf}

\LK Landauers zentrales Modell (Fig.~1 in [B1]) ist ein klassisches Teilchen in
einem bistabilen Potentialtopf:

\begin{itemize}[leftmargin=1.4em,itemsep=0.2em]
  \item linkes Minimum: Zustand $0$;
  \item rechtes Minimum: Zustand $1$;
  \item die Operation \textup{RESTORE TO ONE} zwingt das Teilchen unabhängig vom
        Ausgangszustand in den rechten Topf.
\end{itemize}

Das ist ein klassisches makroskopisches Modell. Es erfasst nicht die volle Breite
moderner Quanteninformationskonzepte, bleibt aber als Analogie für bistabile
Systeme gültig --- Spin-½-Systeme sind formal äquivalent. Entscheidend ist, was
das Modell \emph{nicht} leistet: Es zeigt, unter welchen physikalischen
Bedingungen Dissipation unvermeidlich ist. Es zeigt \emph{nicht}, dass diese
Bedingungen immer dann erfüllt sind, wenn ein Bit im informationstheoretischen
Sinne gelöscht wird.

\begin{center}
\small
\begin{tabular}{@{}p{0.45\textwidth}p{0.45\textwidth}@{}}
\toprule
\textbf{Landauers Modell leistet} & \textbf{Landauers Modell leistet nicht} \\
\midrule
Beweis für physikalische Träger-Operationen & Beweis für informationstheoretisches Löschen \\
untere Schranke für \textup{RESTORE TO ONE} & Aussage über abstrakte Bits \\
Gültigkeit für ein thermalisiertes Ensemble & Gültigkeit für beliebige Rechenprozesse \\
physikalische Entropie $S = k\ln W$ & Brücke zur Shannon-Entropie \\
\bottomrule
\end{tabular}
\end{center}

\subsection{Die Entropieberechnung}

\LK Landauer berechnet die Entropieänderung beim Rücksetzen als
\begin{equation}
  \Delta S \;=\; k \ln 2 - k \ln 1 \;=\; k \ln 2 .
  \label{eq:landauer}
\end{equation}

Diese Rechnung setzt voraus:
\begin{enumerate}[leftmargin=1.6em,itemsep=0.2em]
  \item \textbf{vorher:} Der Träger kann 2 physikalische Zustände annehmen ($W=2$);
  \item \textbf{nachher:} Der Träger kann nur 1 physikalischen Zustand annehmen
        ($W=1$).
\end{enumerate}

Beide Voraussetzungen sind Aussagen über den Träger. Keine der beiden ist eine
Aussage über den Informationsgehalt.

% ===================== 4 =====================
\section{Der Einwand und seine Reichweite}

\LB Die Entropieänderung in Gleichung~\eqref{eq:landauer} ist eine Eigenschaft des
\textbf{physikalischen Trägers}, nicht der Information. Löscht man ein Bit
\emph{interpretativ} --- indem man die Bedeutung umdefiniert, den Wert für
ungültig erklärt, die Adresse aus der Tabelle streicht ---, so ändert sich die
Zahl der physikalischen Zustände des Trägers \textbf{nicht}. Es gilt weiterhin
$W = 2$, also
\begin{equation}
  \Delta S \;=\; k \ln 2 - k \ln 2 \;=\; 0 .
  \label{eq:interpretativ}
\end{equation}

\begin{satz}
Landauers Prinzip greift nur dann, wenn die Operation den Zustandsraum des Trägers
physikalisch verkleinert. Eine rein interpretative Löschung verkleinert ihn nicht.
\end{satz}

\LB \textbf{Beweis.} Sei $W$ die Zahl der physikalischen Mikrozustände des
Trägers. Eine interpretative Löschung erklärt einen Zustand für nicht mehr gültig;
sie ist eine Operation auf der Beschreibungsebene und lässt den physikalischen
Zustandsraum unverändert. Die Entropie des Trägers bleibt $S = k\ln W$. Erst wenn
der Träger selbst auf einen einzigen Zustand gezwungen wird ($W' = 1$), ändert sich
die Entropie auf $S' = 0$. Das ist eine physikalische Operation am Träger, nicht am
Informationsgehalt. \hfill\qedsq

\subsection{Wogegen dieser Einwand nicht gerichtet ist}

\LB Hier ist Genauigkeit nötig, sonst richtet sich der Einwand gegen einen Gegner,
den es nicht gibt. Gleichung~\eqref{eq:interpretativ} widerlegt Landauer
\emph{nicht}. Landauer behauptet an keiner Stelle, dass eine begriffliche
Umwidmung Wärme kostet; niemand in der Fachliteratur behauptet das. Was
Gleichung~\eqref{eq:interpretativ} zeigt, ist ausschließlich, dass die beiden
Vorgänge --- interpretative Löschung und Träger-Rücksetzung --- auseinanderfallen
und dass nur der zweite unter Landauers Schranke fällt.

Das ist keine Kleinigkeit, aber es ist auch nicht das, wonach es auf den ersten
Blick aussieht. Der Ertrag ist ein \emph{Zuständigkeitsbefund}, keine Widerlegung:

\begin{quote}
\LB Landauers Schranke ist zuständig für Träger-Operationen, die den
Zustandsraum verkleinern. Für alles andere, was in der Umgangssprache „Löschen“
heißt, ist sie nicht zuständig --- weder bestätigend noch widerlegend.
\end{quote}

\LB Die eigentliche Frage verschiebt sich damit, und sie ist schwerer als die
ursprüngliche: \emph{Erzwingt informationstheoretisches Löschen eine
zustandsraumverkleinernde Träger-Operation?} Bennett [B7][B8] hat gezeigt, dass es
das im Allgemeinen nicht tut --- jede logisch irreversible Funktion lässt sich
reversibel ausführen, wenn man die Zwischenzustände mitführt. Erst der endgültige
Verzicht auf Träger-Zustände kostet. Damit ist die Antwort für den allgemeinen
Fall negativ, und das ist eine Antwort, die Landauer 1961 noch nicht geben konnte.

% ===================== 5 =====================
\section{Landauers eigene Einschränkung}

\LK Landauer hat eine Einschränkung formuliert, die in der späteren Rezeption
meist übergangen wird:

\begin{quote}
\itshape
„Our reset operation, in the preceding discussion, was applied to a thermal
equilibrium ensemble. In actuality we would like to know what happens in a
particular computing circuit which will work on information which has not yet been
thermalized \ldots“ \upshape [B1]
\end{quote}

Übersetzung: \emph{Unsere Rücksetz-Operation wurde in der vorangegangenen
Diskussion auf ein thermisches Gleichgewichts-Ensemble angewendet. In der
Wirklichkeit möchten wir wissen, was in einer bestimmten Rechenschaltung
geschieht, die mit Information arbeitet, die noch nicht thermalisiert ist.}

\LK Landauer löst die Frage unmittelbar danach selbst auf: Für eine zufällige
Abfolge von Einsen und Nullen setzt er das statistische Ensemble dem Zeitmittel
gleich und gelangt so zu derselben Dissipation pro Rücksetzung wie für das
thermalisierte Ensemble. Für nicht-zufällige Daten --- etwa wenn alle Eingaben
bereits \textup{ONE} sind --- räumt er ausdrücklich ein, dass keine
Entropieänderung stattfindet.

Daraus folgt dreierlei:

\begin{enumerate}[leftmargin=1.6em,itemsep=0.3em]
  \item \LK Landauer ist sich der Grenzen seines Modells bewusst und arbeitet sie
        selbst heraus. Der Vorwurf, er habe überverallgemeinert, trifft die
        Rezeption, nicht ihn.
  \item \LB Sein Ergebnis $kT\ln 2$ gilt streng für das
        Gleichgewichts-Ensemble. Für korrelierte, nicht-zufällige Daten ist
        $kT\ln 2$ pro Träger eine \emph{obere} Schranke; die zutreffende Schranke
        ist $kT\,H$ mit $H$ der Entropie der Quelle in nat. Das ist keine
        Schwächung des Prinzips, sondern seine korrekte Fassung: Die Schranke
        folgt der tatsächlichen Zustandszahl, nicht der Bitbreite.
  \item \LB Vor allem aber: Auch diese Erweiterung bleibt eine Aussage über
        \emph{Träger-Zustände}. Der Übergang zur abstrakten Information wird auch
        hier nicht vollzogen.
\end{enumerate}

% ===================== 6 =====================
\section{Was Landauer beweist --- und was nicht}

Die populäre Kurzfassung lautet:

\begin{quote}
\emph{„Jedes gelöschte Bit erzeugt $kT\ln 2$ Wärme.“}
\end{quote}

\LB Diese Aussage geht über Landauers Beweis hinaus. Bewiesen ist:

\begin{quote}
\emph{Wird ein physikalisches Zwei-Zustands-System, das sich im thermischen
Gleichgewicht befindet, unabhängig von seinem Ausgangszustand in einen definierten
Zustand gezwungen, so muss mindestens $kT\ln 2$ an Wärme abgegeben werden. Diese
Dissipation ist eine Eigenschaft der physikalischen Träger-Operation.}
\end{quote}

\LK Der Schritt, den Landauers Beweis nicht enthält, ist die Brücke von der
Träger-Operation zur abstrakten Information. Landauer verwendet „bit“ durchgehend
in doppelter Bedeutung --- mal für das physikalische Objekt, mal für die
Informationseinheit. Diese Doppeldeutigkeit ist in der Originalarbeit angelegt und
hat die spätere Rezeption begünstigt, in der aus einer Aussage über Träger eine
Aussage über Information wurde.

\LB Dass jedes informationstheoretische Löschen eine zustandsraumverkleinernde
Träger-Operation erfordert, wird vorausgesetzt, nicht hergeleitet --- und ist nach
Bennett im Allgemeinen falsch.

% ===================== 7 =====================
\section{Die moderne Thermodynamik der Information}

\LK Die Thermodynamik der Information mit Rückkopplung [B5][B6] liefert für den
Löschvorgang die verallgemeinerte Schranke
\begin{equation}
  Q_{\min} \;=\; k T\,(\ln 2 - I), \qquad I \in [0,\ \ln 2],
  \label{eq:feedback}
\end{equation}
wobei $I$ die wechselseitige Information zwischen Messgerät und Systemzustand ist,
gemessen in \textbf{nat} und damit dimensionslos. Äquivalent in der Form von
Sagawa und Ueda: $W_{\text{ext}} \le -\Delta F + kT\,I$.

\begin{itemize}[leftmargin=1.4em,itemsep=0.25em]
  \item \LB Weiß der Experimentator, ob der Träger auf $0$ oder $1$ steht
        ($I = \ln 2$), so kann er ihn reversibel zurücksetzen ---
        $Q_{\min} = 0$.
  \item \LB Weiß er es nicht ($I = 0$), so muss er irreversibel zurücksetzen ---
        $Q_{\min} = kT\ln 2$.
\end{itemize}

\LB Die Wärme entsteht also nicht \emph{wegen} des Verlusts von Information,
sondern wegen des mangelnden Wissens über den Trägerzustand. Die Information ist
das Etikett, das auf den Träger geklebt wird; der Preis hängt daran, wie weit das
Etikett den Träger tatsächlich festlegt.

\LB Und zugleich ist die Symmetrie zu sehen: Das Wissen $I$ ist selbst auf einem
Träger gespeichert. Wer es benutzt, um reversibel zurückzusetzen, hat den Preis
nicht vermieden, sondern verschoben --- auf den späteren Zeitpunkt, an dem der
Träger des Wissens seinerseits zurückgesetzt wird. Das ist Bennetts Auflösung des
Maxwellschen Dämons [B8], und sie verhindert, dass
Gleichung~\eqref{eq:feedback} als Umgehung der Schranke gelesen wird.

% ===================== 8 =====================
\section{Die Perpetuierung in der Rezeption}

\LK Die begriffliche Unschärfe wäre folgenlos geblieben, hätte die nachfolgende
Literatur sie korrigiert. Das Gegenteil ist eingetreten. Die Mehrheit der
Veröffentlichungen nach 1961 übernimmt das Prinzip in verallgemeinerter Form, ohne
auf Landauers eigene Einschränkungen zurückzugehen.

\begin{itemize}[leftmargin=1.4em,itemsep=0.3em]
  \item \LK \textbf{Lairez} [B9] stellt fest, dass Kritik an der Verallgemeinerung
        seit Jahrzehnten vorliegt, in der Breite aber keine Wirkung entfaltet hat.
  \item \LK \textbf{Norton} [B4] zeigt, dass das Standardargument auf einer nicht
        haltbaren Gleichsetzung dynamischer und statischer Wahrscheinlichkeiten
        beruht; \textbf{Earman und Norton} [B3] haben die Linie von Szilárd bis
        Landauer bereits 1999 durchgearbeitet.
  \item \LK \textbf{Hemmo und Shenker} [B10] weisen nach, dass das Prinzip im
        Licht des Multiple-Computations-Theorems nicht universell gültig sein kann.
\end{itemize}

\LSk Der Grund für die Persistenz liegt nicht in einer Widerlegung dieser Kritik,
sondern in strukturellen Faktoren der Wissenschaftskommunikation: Die Formel
$kT\ln 2$ ist griffig und leicht zitierbar. Die Originalarbeit von 1961 wird selten
gelesen; zitiert werden Sekundärquellen, die ihrerseits Sekundärquellen zitieren.
Autoritätsvertrauen gegenüber einem renommierten Industrieforscher tut ein
Übriges. Die Folge ist, dass ein nicht vollzogener Beweisschritt als
„fundamentales Naturgesetz“ gehandelt wird --- eine Kennzeichnung, die in der
Originalarbeit nicht auftaucht.

\LB Diese Diagnose ist eine Plausibilitätsskizze und als solche markiert. Sie
benennt Mechanismen, die einzeln beobachtbar, in ihrem Zusammenwirken aber nicht
quantifiziert sind. Eine belastbare Fassung wäre eine Zitationsanalyse: Anteil der
Arbeiten, die [B1] zitieren und dabei die Ensemble-Einschränkung nennen. Diese
Analyse liegt hier nicht vor.

% ===================== 9 =====================
\section{Die Unschärfe in der Sprache der Digitaltechnik}

\LB Das Problem ist nicht auf Landauer beschränkt. Die gesamte Digitaltechnik
verwendet eine Sprache, die systematisch Träger-Eigenschaften als
Informations-Eigenschaften formuliert. Die Praxis ist so tief verankert, dass sie
von Praktikern selten als problematisch wahrgenommen wird.

\begin{center}
\small
\begin{longtable}{@{}p{0.26\textwidth}p{0.34\textwidth}p{0.30\textwidth}@{}}
\toprule
\textbf{Sprachüblicher Ausdruck} & \textbf{Physikalische Realität (Träger)} & \textbf{Problematische Implikation} \\
\midrule
\endfirsthead
\toprule
\textbf{Sprachüblicher Ausdruck} & \textbf{Physikalische Realität (Träger)} & \textbf{Problematische Implikation} \\
\midrule
\endhead
„ein Bit setzen“ & Spannung an einen Kondensator anlegen & Information habe einen Zustand \\
„ein Bit löschen“ & Träger auf Referenzspannung setzen & Information könne vernichtet werden \\
„ein Bit übertragen“ & Signal auf einer neuen Leitung rekonstruieren & Information sei transportierbar \\
„ein Bit kopieren“ & Spannung auslesen, anderswo erzeugen & Information vervielfältige sich \\
„das RAM hat 8\,GB“ & 8 Milliarden physikalische Speicherzellen & Information sei eine Mengengröße \\
„Bitrate“ & Spannungswechsel pro Sekunde & Informationsmenge pro Zeit \\
\bottomrule
\end{longtable}
\end{center}

In allen Fällen wird eine Eigenschaft des Trägers so formuliert, als wäre sie eine
Eigenschaft der Information. Ingenieure wissen in der Praxis, dass sie mit
Spannungen, Strömen und Ladungen hantieren. In der Beschreibungssprache aber wird
das Bit zum handelnden Subjekt: Es „fließt durch den Bus“, „wird gespeichert“,
„wird gelöscht“. Diese Sprache ist für viele Zwecke praktisch --- sie verführt aber
dazu, physikalische Gesetze unmittelbar auf die Abstraktion anzuwenden, wie es die
populäre Landauer-Rezeption getan hat.

\LB Ein präziserer Sprachgebrauch würde unterscheiden:

\begin{itemize}[leftmargin=1.4em,itemsep=0.25em]
  \item \textbf{Nicht:} „Das Bit hat Energie.“ \quad
        \textbf{Sondern:} „Der Träger dieses Bits hat Energie.“
  \item \textbf{Nicht:} „Das Bit wird gelöscht.“ \quad
        \textbf{Sondern:} „Der Träger wird auf einen definierten Zustand gesetzt.“
  \item \textbf{Nicht:} „Das Bit kostet $kT\ln 2$.“ \quad
        \textbf{Sondern:} „Diese Träger-Operation kostet unter diesen Bedingungen
        mindestens $kT\ln 2$.“
\end{itemize}

Die Verbindung zu Landauer ist direkt: Er verwendet „bit“ in eben dieser
doppeldeutigen Weise. Die Doppeldeutigkeit ist nicht seine Erfindung; sie ist in
der Disziplin angelegt und hat die Verallgemeinerung getragen.

% ===================== 10 =====================
\section{Erkenntnistheoretischer Hintergrund: Hypostasierung}

\LSk Der beschriebene Mechanismus --- eine Abstraktion wird als physikalische
Größe behandelt und mit physikalischen Eigenschaften ausgestattet --- hat einen
Namen: \emph{Hypostasierung} (oder Reifikation), die Verwandlung eines abstrakten
Begriffs in ein vermeintlich reales, substanzielles Objekt.

\LK Das Bit ist eine Abstraktion. Shannon [B11] definierte es als statistische
Eigenschaft von Signalen, nicht als physikalisches Objekt. Wiener formulierte
denselben Sachverhalt ausdrücklich: Information sei Information, nicht Materie
oder Energie [B12]. Diese Klarheit ging in der ingenieurwissenschaftlichen Praxis
verloren. Das Bit wurde zum Objekt, dem man Energie, Entropie und Lokalisation
zuschreiben konnte.

\LSk Der Mechanismus ist in der Wissenschaftsgeschichte nicht neu. Die Entropie,
das elektrische Feld, das Gen --- alle diese Abstraktionen wurden zu bestimmten
Zeitpunkten als substanzielle Dinge behandelt, mit Folgen für die Theorie. Im Fall
der Information führt die Hypostasierung zu einer falschen Physik: Hätte das Bit
selbst Entropie, dann müsste sein Löschen Wärme erzeugen --- unabhängig vom
Träger. Dieser Schluss ist nur dann zwingend, wenn man die Abstraktion mit der
physikalischen Realität verwechselt hat.

\LB Der Abschnitt ist als Skizze markiert, weil er eine Deutung anbietet und
keinen Nachweis führt. Die historischen Parallelen sind Analogien; sie stützen die
These, sie beweisen sie nicht.

% ===================== 11 =====================
\section{Schlussfolgerung}

\begin{enumerate}[leftmargin=1.6em,itemsep=0.35em]
  \item \LB Landauers Beweis ist für die physikalische Operation
        \textup{RESTORE TO ONE} an einem thermischen Gleichgewichts-Ensemble
        korrekt und präzise.
  \item \LB Er zeigt eine untere Schranke für \emph{Träger-Operationen} --- keine
        universelle Aussage über abstrakte Information.
  \item \LB Die Übertragung auf den informationstheoretischen Begriff des Bits
        setzt voraus, dass jedes informationstheoretische Löschen mit einer
        zustandsraumverkleinernden Träger-Operation verbunden ist. Diesen
        Zusammenhang beweist Landauer nicht --- und nach Bennett gilt er im
        Allgemeinen nicht.
  \item \LK Die Doppeldeutigkeit von „bit“ ist in der Originalarbeit angelegt und
        hat die Verallgemeinerung in der Rezeption begünstigt.
  \item \LK Landauer selbst diskutiert die Grenzen seines Ensemble-Arguments und
        räumt ein, dass nicht-zufällige Eingabedaten die Dissipation unter
        $kT\ln 2$ senken.
  \item \LB Was hier nicht gezeigt wird: dass das Landauer-Prinzip falsch sei.
        Gezeigt wird, dass es für einen anderen Gegenstand zuständig ist als den,
        für den es meist zitiert wird.
\end{enumerate}

% ===================== 12 =====================
\section{Ausblick}

\LB Eine Formulierung, die den tatsächlichen Beweis wiedergibt, lautet:

\begin{quote}
\emph{Die physikalische Operation, die ein bistabiles System im thermischen
Gleichgewicht unabhängig von seinem Ausgangszustand in einen definierten Zustand
versetzt, erfordert mindestens $kT\ln 2$ an Dissipation. Dient dieses System als
Informationsträger und entspricht die Operation dem Löschen eines Bits, so gilt
die Schranke für die zugehörige Träger-Operation --- nicht für das
informationstheoretische Löschen als solches.}
\end{quote}

\LH \textbf{Offene Frage.} Unter welchen Bedingungen erzwingt ein gegebener
Rechenprozess eine zustandsraumverkleinernde Träger-Operation? Bennett beantwortet
den Grenzfall (im Prinzip nie, wenn man alles mitführt); die praktisch
interessante Frage ist die quantitative: Wie viel Träger-Zustandsraum muss eine
Architektur bei gegebener Speichertiefe endgültig aufgeben? Das ist dieselbe
offene Frage wie der kontinuierliche Entropiestrom in A280 --- von der anderen
Seite gestellt.

% ===================== 13 =====================
\section{Schichtstatus-Zusammenfassung}

\begin{center}
\small
\begin{longtable}{@{}p{0.72\textwidth}l@{}}
\toprule
\textbf{Aussage} & \textbf{Status} \\
\midrule
\endfirsthead
\toprule
\textbf{Aussage} & \textbf{Status} \\
\midrule
\endhead
Definition Information (abstrakt, ohne Energie) & \LS \\
Definition Träger (physikalisch, mit Energie) & \LS \\
Wörterbuch-Zuordnung A280 $\leftrightarrow$ A281 & \LS \\
Satz 1: Energie hängt am Träger, nicht am Bit & \LB \\
Bus-Beispiel: Busbreite ist Hardware, nicht Information & \LB \\
Landauers Doppeltopf-Modell und \textup{RESTORE TO ONE} & \LK \\
$\Delta S = k\ln 2$ setzt $W{=}2 \to W{=}1$ am Träger voraus & \LK \\
Satz 2: interpretative Löschung verkleinert $W$ nicht & \LB \\
Der Einwand ist ein Zuständigkeitsbefund, keine Widerlegung & \LB \\
Bennett: logisch irreversible Funktionen reversibel ausführbar & \LB \\
Landauers Ensemble-Einschränkung (Originalzitat) & \LK \\
Landauers eigene Auflösung für zufällige Daten & \LK \\
Für korrelierte Daten gilt $kT\,H$, nicht $kT\ln 2$ & \LB \\
Landauer verwendet „bit“ doppeldeutig & \LK \\
Rückkopplungsschranke $Q_{\min} = kT(\ln 2 - I)$, $I$ in nat & \LK \\
Das Wissen $I$ ist selbst trägergebunden (Bennett) & \LB \\
Rezeptionsbefunde Lairez / Norton / Hemmo-Shenker & \LK \\
Strukturelle Gründe der Persistenz & \LSk \\
Sprachtabelle Träger vs.\ Information & \LB \\
Hypostasierung als Deutungsrahmen & \LSk \\
Historische Parallelen (Entropie, Feld, Gen) & \LSk \\
Kein Nachweis der Falschheit des Landauer-Prinzips & \LB \\
Quantitative Fassung des unvermeidbaren Zustandsraumverzichts & \LH \\
\bottomrule
\end{longtable}
\end{center}

\textbf{Bilanz:} 10$\times$ \LB $\cdot$ 8$\times$ \LK $\cdot$ 3$\times$ \LS
$\cdot$ 3$\times$ \LSk $\cdot$ 1$\times$ \LH $=$ 25 Einträge.

% ===================== ANHANG =====================
% ===================== BELEGE =====================
\clearpage
\phantomsection
\addcontentsline{toc}{section}{Belege}
\section*{Belege}

\begingroup
\small
\begin{description}[leftmargin=3.2em,style=sameline,itemsep=0.3em]

\item[{[B1]}] R. Landauer, \emph{Irreversibility and Heat Generation in the
Computing Process}, IBM J.\ Res.\ Dev.\ \textbf{5}, 183 (1961).

\item[{[B2]}] C.~H. Bennett, \emph{Logical Reversibility of Computation}, IBM J.\
Res.\ Dev.\ \textbf{17}, 525 (1973).

\item[{[B3]}] J. Earman und J.~D. Norton, \emph{Exorcist XIV: The Wrath of
Maxwell's Demon. Part II: From Szilard to Landauer and Beyond}, Stud.\ Hist.\
Phil.\ Mod.\ Phys.\ \textbf{30}, 1 (1999).

\item[{[B4]}] J.~D. Norton, \emph{All Shook Up: Fluctuations, Maxwell's Demon and
the Thermodynamics of Computation}, Entropy \textbf{15}, 4432 (2013).

\item[{[B5]}] T. Sagawa und M. Ueda, \emph{Nonequilibrium Thermodynamics of
Feedback Control}, Phys.\ Rev.\ E \textbf{85}, 021104 (2012).

\item[{[B6]}] J.~M.~R. Parrondo, J.~M. Horowitz und T. Sagawa,
\emph{Thermodynamics of Information}, Nature Physics \textbf{11}, 131 (2015).

\item[{[B7]}] J. Bub, \emph{Maxwell's Demon and the Thermodynamics of
Computation}, Stud.\ Hist.\ Phil.\ Mod.\ Phys.\ \textbf{32}, 569 (2001).

\item[{[B8]}] C.~H. Bennett, \emph{The Thermodynamics of Computation --- a
Review}, Int.\ J.\ Theor.\ Phys.\ \textbf{21}, 905 (1982).

\item[{[B9]}] \emph{(unvollständig --- Kennung nachzutragen)} D. Lairez,
\emph{Landauer's Principle: A Critical Assessment}, arXiv-Preprint (2024).

\item[{[B10]}] M. Hemmo und O. Shenker, \emph{The Multiple-Computations Theorem
and the Physics of Singling Out a Computation}, The Monist \textbf{105}, 175
(2022).

\item[{[B11]}] C.~E. Shannon, \emph{A Mathematical Theory of Communication}, Bell
System Technical Journal \textbf{27}, 379 (1948).

\item[{[B12]}] \emph{(rekonstruiert --- Seitenangabe zu prüfen)} N. Wiener,
\emph{Cybernetics: or Control and Communication in the Animal and the Machine},
MIT Press, Cambridge, MA (1948).

\end{description}
\endgroup

\vspace{1em}
\noindent
{\small\itshape Einträge, die als \emph{unvollständig} oder \emph{rekonstruiert}
markiert sind, wurden aus dem Zitationskontext ergänzt bzw.\ tragen einen
offenen Platzhalter; sie sind vor der Veröffentlichung gegen die
Arbeitsunterlagen zu prüfen.}

\end{document}
